\subsubsection{Shortest distance, diameter and average path length}

In a Watts-Strogatz small world graph, edges connect two arbitrary nodes $i$ and $j$ through a path, a sequence of edges. A graph is connected if any $i$ can be reached from any $j$ through any path of any length. With the assumption that every graph is connected, define the diameter and average path length. Begin with dist($i, j$), which is the shortest-path distance between $i$ and $j$, i.e. the least number of edges which connect the two nodes. The largest possible shortest path between two nodes in the graph is the diameter, written as 
\begin{align}
    \text{diam(G)} = \max_{i,j} \text{dist}(i,j).
    \label{eq:WSdiameter}
\end{align}

The average path length is instead the average shortest path length between $i$ and $j$ and is defined as follows:
\begin{align}
    \ell(\text{G}) = \frac{1}{N(N-1)}\sum_{i\neq j}\text{dist}(i,j)
    \label{eq:WSavgshortestpath}
\end{align}

A graph can also be disconnected if there are nodes or sets of nodes that cannot be reached from any other node. The diameter and average path length would still be defined in the same way as for a connected graph, but the numerical value would be different once computed.
